\appsection{Resumé}

Táto kapitola zhŕňa obsah dizertačnej práce v slovenskom jazyku a je určená čitateľom, ktorí nepracujú s anglickým plným textom. Práca sa zaoberá Internetom hodnoty (Internet of Value) a navrhuje architektúru pre agentové rozhodovanie nad heterogénnymi digitálnymi aktívami.

\subsection{Úvod a motivácia}

Internet hodnoty je prostredie, v ktorom je možné digitálne reprezentovať, prenášať a vyhodnocovať rôznorodé aktíva. Patria sem nezameniteľné tokeny, produktívne aktíva ako stakované deriváty, hlasovacie a správcovské práva v decentralizovaných organizáciách, finančné nástroje aj prístupové oprávnenia. Blockchain ako infraštruktúra poskytuje vyrovnanie transakcií, nemennosť záznamov a programovateľnú logiku prostredníctvom inteligentných kontraktov. To, čo blockchain neposkytuje, je vrstva, ktorá by pred samotným vykonaním akcie umožnila posúdiť, akú hodnotu má dané aktívum pre konkrétny zámer používateľa alebo decentralizovanej organizácie.

V praxi to znamená, že aj keď je možné aktívum kúpiť, predať alebo presunúť na úrovni protokolu, rozhodnutie o tom, či má daná transakcia v danom kontexte zmysel, ostáva mimo reťazca a väčšinou bez zachytenej argumentácie. Súčasné riešenia tento problém zužujú na trhovú cenu, hoci hodnota mnohých aktív vyplýva skôr z úžitku, výnosu, správcovskej váhy alebo prístupových práv. Práca preto skúma, ako sformalizovať zámer používateľa ako vstup do hodnotenia a ako naviazať následné akcie na overiteľnú stopu rozhodnutia.

Druhým motivačným bodom je rastúca úloha autonómnych a poloautonómnych agentov. Veľké jazykové modely a štruktúrované rozhrania pre orchestráciu nástrojov umožňujú delegovať časť získavania informácií a štruktúrovanej analýzy na softvérových agentov. Súčasne však platí, že agent nedokáže poskytnúť tvrdé záruky o správnosti svojho výstupu. Práca preto skúma, za akých podmienok môže byť agentové odporúčanie použité ako vstup do rozhodnutia o aktívach a aké hranice musia byť vynútené mimo agenta, aby výsledné rozhodnutie zostalo overiteľné a obmedzené politikou.

\subsection{Teoretické východiská}

Druhá kapitola práce zhŕňa základné pojmy potrebné na pochopenie ďalšieho výkladu. Vysvetľuje princíp blockchainu ako distribuovanej databázy, rolu konsenzuálnych mechanizmov a rozdiel medzi verejnými a povolenými sieťami. Inteligentné kontrakty sú predstavené ako programy, ktoré sa vykonávajú deterministicky na sieti uzlov a ktorých stav je verejne overiteľný. Práca opisuje aj relevantné štandardy tokenov, najmä ERC-20 pre zameniteľné tokeny a ERC-721 a ERC-1155 pre nezameniteľné a kombinované aktíva.

Osobitná pozornosť je venovaná abstrakcii účtu, ktorá umožňuje vytvárať programovateľné účty so smart kontraktovou logikou namiesto klasických externe vlastnených adries. K týmto témam patria aj peňaženky pre agentov, delegácia podpisu, modulárne podpisové schémy a obmedzenia, ktoré je možné vynútiť na úrovni účtu. Tieto stavebné prvky tvoria základ, na ktorom práca neskôr stavia svoju vrstvenú architektúru a rozhodovaciu logiku.

\subsection{Analýza súčasného stavu}

Tretia kapitola sa venuje existujúcim prístupom v oblastiach, ktoré sa s prácou tematicky prekrývajú. V oblasti vyjadrenia zámeru sa rozoberajú intentové protokoly a riešiteľské siete (solvers), ktoré umožňujú používateľovi vyjadriť žiadaný výsledok bez špecifikácie konkrétnej cesty jeho dosiahnutia. Analyzované sú aj systémy delegovaného vykonávania, kde transakciu zostavuje a podáva iná strana ako vlastník aktíva, a riešenia pre tzv. policy wallets, ktoré obmedzujú typy transakcií na úrovni účtu.

V oblasti agentových systémov sa práca venuje rámcom pre orchestráciu jazykových modelov, protokolu Model Context Protocol, agentovým identitám a reputačným mechanizmom. V oblasti správy decentralizovaných organizácií sú zhrnuté nástroje pre on-chain hlasovanie, správu pokladnice a nástroje pre podporu rozhodovania, ako sú napríklad systémy typu Polkadot OpenGov. V oblasti hodnotenia aktív sa analyzujú metódy oceňovania nezameniteľných tokenov, oceňovania produktívnych aktív cez výnosové modely a viackriteriálne metódy rozhodovania.

Z analýzy vyplývajú tri medzery v súčasnom stave. Prvá medzera je v hodnotení, kde sa hodnota redukuje na trhovú cenu aj v prípadoch, keď to je nevhodné. Druhá medzera je v účtovom modeli, ktorý je orientovaný na autentifikáciu a delegáciu, no nepokrýva vynucovanie politík viazaných na zámer a overiteľný záznam o rozhodnutí. Tretia medzera je v správe pokladníc, kde objem návrhov narastá rýchlejšie, než ich dokážu kvalifikovane posúdiť ľudskí hodnotitelia. Tieto medzery vytvárajú priestor pre rozhodovaciu vrstvu, ktorá prepája zámer, dôkazy, politiku a vykonanie.

\subsection{Ciele dizertačnej práce}

Cieľom práce je navrhnúť a overiť architektúru, ktorá umožňuje agentové rozhodovanie nad rôznorodými aktívami v prostredí Internetu hodnoty tak, aby bolo zachované hodnotenie podľa zámeru, dohľadateľnosť dôkazov, schválenie človekom a obmedzenia politiky na úrovni účtu. Hlavným prínosom je formálny model Intent-Conditioned Value Projection (ICVP), v ktorom zámer používateľa alebo organizácie pôsobí ako projekcia nad priestorom heterogénnych dôkazov. Zámer neurčuje len dôležitosť dimenzií dostupnosti, trvácnosti a úžitku, ale aj to, čo tieto dimenzie pre dané aktívum a danú akciu znamenajú.

Práca formuluje tri výskumné otázky:

\begin{enumerate}
   \item \textbf{VO1.} Aký formálny model dokáže reprezentovať zámer používateľa alebo organizácie ako projekciu nad heterogénnymi dôkazovými priestormi Internetu hodnoty a produkovať inštanciované hodnotové dimenzie a overiteľné rozhodovacie signály pre aktíva a akcie nad aktívami?
   \item \textbf{VO2.} Akou metódou je možné Intent-Conditioned Value Projection uviesť do prevádzky v agentovom systéme podpory rozhodovania so zachovaním dohľadateľnosti dôkazov, schválenia človekom a politicky obmedzenej prípravy transakcií?
   \item \textbf{VO3.} Do akej miery Intent-Conditioned Value Projection zlepšuje kvalitu rozhodnutia, relevantnosť dôkazov a auditovateľnosť pri rozhodnutiach o nezameniteľných tokenoch, produktívnych aktívach a správe pokladnice v porovnaní s jednorozmerným variantom a variantom s pevnými váhami?
\end{enumerate}

Okrem výskumných otázok práca vymedzuje aj dve bezpečnostné vlastnosti, a to úplnosť dôkazov pri kladnom rozhodnutí a viazanosť autorizácie transakcie na platnú stopu rozhodnutia.

\subsection{Návrh riešenia}

Návrh je najrozsiahlejšia časť práce. Architektúra má štyri vrstvy. Lokálna vrstva zodpovedá za interakciu s používateľom, za prijatie zámeru a za schválenie navrhnutých akcií. Dátová a dôkazová vrstva zhromažďuje informácie potrebné na hodnotenie. Agentová vrstva vykonáva samotné rozhodovanie a prípravu transakcie. On-chain vrstva vykonáva schválené akcie a vynucuje politiku na úrovni účtu.

Jadrom modelu je formalizácia zámeru ako projekcie. Používateľ alebo organizácia vyjadrí, čo chce v danej situácii dosiahnuť. Zámer sa premietne do dôkazového priestoru a vyberá z neho len tie atribúty aktíva, ktoré sú pre danú akciu relevantné. Na základe tohto výberu sa vytvoria tri dimenzie hodnotenia, ktoré model označuje ako dostupnosť, trvácnosť a úžitok. V kontexte nezameniteľného tokenu môže dostupnosť znamenať likviditu na trhu, trvácnosť stabilitu kolekcie a úžitok prístup k danej komunite alebo službe. V kontexte stakovaného aktíva môže dostupnosť znamenať možnosť okamžitej výplaty, trvácnosť bezpečnosť validátorov a úžitok očakávaný výnos. V kontexte správcovského návrhu môže dostupnosť znamenať pripravenosť pokladnice, trvácnosť dlhodobý dopad návrhu a úžitok prínos pre verejné statky.

Aby bolo možné rozhodnutie overiť, model zavádza pojem požadovaného dôkazu (RequiredEvidence). Pre danú triedu aktíva a daný zámer je vopred určené, aké typy dôkazov musia byť k dispozícii, aby bolo možné rozhodnúť kladne. Ak dôkazy chýbajú alebo nie sú overiteľne získané, rozhodnutie nemôže byť kladné a systém odporúča rozhodnutie odložiť alebo zamietnuť. Takto je definovaná vlastnosť úplnosti dôkazov pri kladnom rozhodnutí.

Výsledkom hodnotenia je stopa rozhodnutia (valuation trace). Ide o štruktúrovaný artefakt, ktorý zaznamenáva zadaný zámer, identifikátor aktíva, vybraté dôkazy aj s ich pôvodom, hodnoty inštanciovaných dimenzií, finálne rozhodnutie z množiny pokračovať, odložiť alebo zamietnuť, a metadáta potrebné na audit. Stopa nie je voľný textový popis, ale validovateľný objekt, voči ktorému je možné spustiť deterministické kontroly. Stopu je možné archivovať mimo reťazca aj v komprimovanej forme na reťazci ako záväzok, takže neskorší audit dokáže overiť, že rozhodnutie skutočne vychádzalo z konkrétnych dôkazov v konkrétnom čase. Tento prístup oddeluje rozsiahle dôkazové dáta od relatívne malého záväzku, ktorý je potrebné držať v reťazci.

Schémy požadovaných dôkazov sú navrhnuté tak, aby boli rozšíriteľné. Pre novú triedu aktív stačí pridať definíciu, aké atribúty sú pre daný zámer povinné, aké voliteľné a aké zdroje sa považujú za prijateľné. Validátor pracuje so schémou ako s deklaratívnym popisom a nevyžaduje úpravu vlastnej logiky. To umožňuje, aby organizácia postupne pokrývala nové typy rozhodnutí bez zásahu do jadra systému. Práca uvádza schémy pre tri skúmané triedy a opisuje, ako boli iteratívne dolaďované na základe expertnej spätnej väzby.

Agentová vrstva pracuje s viacerými spolupracujúcimi agentmi. Jeden agent vyhľadáva a získava dôkazy zo zdrojov pripojených cez Model Context Protocol. Druhý agent na základe zámeru a dôkazov navrhuje hodnoty dimenzií a finálne rozhodnutie. Tretí agent zostavuje návrh transakcie tak, aby zodpovedal stope rozhodnutia a obmedzeniam politiky účtu. Validátor mimo agentovej vrstvy potom deterministicky kontroluje, či je stopa kompletná, či sú dôkazy v poriadku a či navrhnutá transakcia zodpovedá rozhodnutiu. Až po úspešnej validácii sa návrh predloží človeku na schválenie.

On-chain vrstva je založená na primitíve, ktoré práca označuje ako NFT ako účet (Non-Fungible Token as Account, NFTAA). Aktíva alebo skupiny aktív sú spravované cez programovateľný účet reprezentovaný nezameniteľným tokenom. Účet má pripojenú politiku, ktorá definuje, aké typy transakcií sú povolené, kto ich môže schvaľovať a aké stopy rozhodnutia sú prijímané ako podklad. Účet pri každej akcii kontroluje, či existuje platná stopa, či bola schválená ľudským potvrdením a či transakcia spadá do povolenej politiky. Takto je definovaná druhá bezpečnostná vlastnosť, teda viazanosť autorizácie na stopu.

Keďže účet je sám nezameniteľný token, je možné ho prenášať, dediť, prenajímať alebo zveriť do správy bez toho, aby sa stratila pripojená politika a história rozhodnutí. Pre organizáciu to znamená, že spravovanie pokladnice môže byť rozdelené medzi viacero účtov s rôznymi politikami, pričom každý účet má jasne vymedzený rozsah pôsobnosti. Pre používateľa to znamená, že jeden účet môže obsahovať skupinu aktív viazaných na konkrétny cieľ, napríklad dlhodobé sporenie, krátkodobé obchodovanie alebo účasť v správe konkrétneho protokolu. Tieto kontexty sa medzi sebou nemiešajú a každý má vlastnú politiku schvaľovania.

Súčasťou návrhu je aj model hrozieb. Práca explicitne pomenúva, čo systém nesľubuje. Nesľubuje, že jazykový model je vždy správny, že experti predstavujú absolútnu pravdu, ani že agent má autonómne kontrolovať aktíva používateľa. Sľubuje, že kladné rozhodnutie je možné len pri prítomnosti požadovaných dôkazov, že transakcia môže byť autorizovaná len pri platnej stope a že politika účtu obmedzuje rozsah povolených akcií. Toto rozdelenie zodpovednosti medzi agenta, validátora, človeka a on-chain vrstvu je zámerné a tvorí základ bezpečnostných vlastností práce.

\subsection{Experimentálne overenie a hodnotenie}

Hodnotenie nie je založené na jednej metrike. Kombinuje deterministickú validáciu stopy, expertný referenčný protokol a porovnanie troch variantov modelu. Plný variant zodpovedá navrhnutej projekcii zámeru s úplnými dimenziami. Prvý zredukovaný variant používa len jednu dominantnú dimenziu rozhodnutia, čo zodpovedá rozšírenej praxi redukcie hodnoty na trhovú cenu. Druhý zredukovaný variant ponecháva všetky tri dimenzie, ale s pevnými váhami a bez projekcie zámeru, čím zachytáva prípady, keď organizácia používa fixnú rubrikovú metodiku.

Hodnotenie prebehlo v troch režimoch, ktoré pokrývajú rozdielne triedy aktív. Prvý režim sa týka kolekcie Pudgy Penguins a reprezentuje rozhodnutia o nezameniteľných tokenoch, kde sa miešajú trhové, komunitné a prístupové aspekty. Druhý režim sa týka stakovaného Etheru (stETH) a reprezentuje rozhodnutia o produktívnych aktívach s výnosom, výplatným profilom a rizikom protokolu. Tretí režim sa týka správcovských návrhov v sieti Polkadot OpenGov a reprezentuje rozhodnutia o pridelení prostriedkov pokladnice. Spolu pokrývajú dvanásť úloh, ktoré boli pre každý variant modelu spustené nezávisle.

Expertný protokol slúži ako referenčný bod. Traja experti nezávisle vyhodnotili tie isté úlohy a poskytli rozhodnutie, dôvody a odhad istoty. Ich rozhodnutia neslúžia ako absolútna pravda, ale ako referencia, voči ktorej sa hodnotí relevantnosť dôkazov, súlad rozhodnutia a kvalita argumentácie. Práca pri tom dôsledne odlišuje, čo môže experimentálne overenie dokázať a čo nie. Overenie netvrdí, že navrhnutý systém je univerzálne lepší ako človek. Tvrdí, že plný variant modelu zachováva viac relevantných dôkazov, produkuje stopy, ktoré prejdú validáciou, a v skúmaných režimoch dosahuje rozhodnutia bližšie k expertnej referencii než zredukované varianty.

Metriky hodnotenia pokrývajú niekoľko rovín. Na rovine rozhodnutia sa sleduje zhoda finálneho výstupu s expertnou referenciou. Na rovine dôkazov sa sleduje, aký podiel expertmi označených relevantných dôkazov skutočne figuruje v stope, a naopak, aký podiel získaných dôkazov experti označili za podstatné. Na rovine argumentácie sa sleduje, či zdôvodnenie vo vyhotovenej stope odkazuje na konkrétne dôkazy a či neobsahuje tvrdenia bez podkladu. Na rovine validácie sa sleduje, aký podiel vygenerovaných stôp prejde deterministickou kontrolou bez zásahu človeka. Tieto roviny sa interpretujú spolu, pretože samostatne by mohli vytvárať skreslený obraz.

Súčasťou hodnotenia je aj analýza robustnosti a nákladov. Robustnosť bola testovaná zmenou formulácie zámeru pri zachovaní jeho významu a sledovaním stability rozhodnutia. Náklady boli merané jednak v podobe spotrebovaných tokenov modelu pri behu agentov, jednak v podobe gas nákladov pri vykonaní transakcií cez NFTAA. Z výsledkov vyplýva, že réžia vrstvy rozhodovania je zvládnuteľná v kontexte typických aktívnych operácií a že najväčší podiel nákladov nesie získavanie a validácia dôkazov, nie samotná on-chain logika. Práca tieto zistenia interpretuje opatrne a vyhýba sa zovšeobecneniam mimo skúmaných režimov.

Z porovnania troch variantov modelu vyplývajú konzistentné zistenia naprieč režimami. Variant s jednou dimenziou dosahuje rozumné výsledky tam, kde je rozhodnutie skutočne dominantne určené jedným faktorom, napríklad pri jednoduchom predaji za trhovú cenu. V zložitejších rozhodnutiach, kde sa miešajú úžitok, riziko a dlhodobý dopad, však tento variant stráca informáciu a častejšie sa odchyľuje od expertnej referencie. Variant s pevnými váhami funguje lepšie pri rutinných úlohách v rámci jednej triedy aktív, no zlyháva pri zmene zámeru, keď by tie isté dôkazy mali mať iný význam. Plný variant si zachováva schopnosť reagovať na zmenu zámeru a v skúmaných úlohách najlepšie kombinuje pokrytie dôkazov a kvalitu argumentácie.

Reprodukovateľnosť hodnotenia zabezpečujú priložené artefakty. Patria k nim schémy požadovaných dôkazov pre jednotlivé triedy aktív, validátor stopy, datasety úloh a expertné anotácie. Tieto artefakty umožňujú zopakovať hodnotenie aj porovnanie s budúcimi variantmi modelu alebo s inými agentovými systémami.

\subsection{Záver}

Práca prepája hodnotenie podľa zámeru, agentové rozhodovanie, overiteľnú stopu a obmedzené vykonanie na úrovni účtu. Model Intent-Conditioned Value Projection chápe zámer ako projekciu nad dôkazmi a inštanciuje dostupnosť, trvácnosť a úžitok podľa triedy aktíva a kontextu akcie. Architektúra prevádza tento model do agentového systému a primitívum NFT ako účet viaže on-chain autorizáciu na platnú stopu rozhodnutia. Hodnotenie v troch režimoch potvrdzuje prínos plného variantu oproti zredukovaným verziám.

Práca má niekoľko obmedzení. Hodnotenie pokrýva tri triedy aktív a vychádza z konkrétnych zdrojov dôkazov, takže závery sa nedajú bez ďalšieho overenia preniesť na ľubovoľné aktívum alebo doménu. Spoľahlivosť agentov závisí od kvality jazykového modelu a od stability nástrojov, ktoré sa v čase menia. Politiky účtu pokrývajú typické prípady, no v reálnej prevádzke ich treba doladiť pre konkrétnu organizáciu. Expertná referencia vychádza z troch nezávislých posudzovateľov, čo zachytáva bežnú rozptylovosť názoru, ale nenahrádza dlhodobé pozorovanie v produkčnej prevádzke.

Budúca práca môže pokračovať vo viacerých smeroch. Patrí k nim rozšírenie schém požadovaných dôkazov na ďalšie triedy aktív, integrácia s ďalšími on-chain prostrediami, hlbšie skúmanie ľudského schvaľovania pri väčšom objeme rozhodnutí a štúdium dlhodobého správania systému pri zmenách trhových podmienok. Zaujímavým smerom je aj skúmanie agregácie viacerých nezávislých agentov, ktorých výstupy by sa porovnávali pred tým, ako vznikne finálna stopa, čo by mohlo zlepšiť robustnosť rozhodnutia v hraničných prípadoch. Ďalším smerom je prepojenie navrhnutej rozhodovacej vrstvy s formálnou verifikáciou politík účtu, aby bolo možné dokázať vlastnosti politiky ešte pred jej nasadením.

Hlavným metodologickým príspevkom práce je rozdelenie zodpovednosti medzi softvérového agenta, validátor, človeka a on-chain vrstvu. Agent navrhuje, validátor kontroluje, človek schvaľuje a účet vynucuje. Tento prístup môže slúžiť ako východisko pre ďalší výskum aj pre praktické nasadenie v decentralizovaných organizáciách a používateľských peňaženkách.
